% SEGMENT OLP-0562-S01
\documentclass[../../../include/open-logic-section]{subfiles}

\begin{document}

\olfileid{sth}{spine}{foundation}

\olsection{அடித்தளம்}

நாம் \emph{கிட்டத்தட்ட} மட்டுமே முடித்திருக்கிறோம்---\emph{முழுமையாக}
முடிக்கவில்லை---ஏனெனில் \emph{ஒவ்வொரு} கணமும் ஏதோ ஒரு $V_\alpha$ இல் உள்ளது,
அதாவது ஒவ்வொரு கணமும் ஏதோ ஒரு படிநிலையில் உருவாகிறது, என்று
$\ZFminus$ இலுள்ள எதுவும் உறுதிசெய்வதில்லை.

படிநிலைமுறைக்கு வெளியே கணங்கள் உள்ளனவா என்பது குறித்து நாம்
\emph{கவலைப்படத் தேவையில்லை} என்று கூறுவதற்கு ஓரளவு நேரடியான (கணித)
பொருள் உள்ளது. (அங்கு ஏதேனும் இருந்தால், அவற்றைப் புறக்கணிக்கலாம்.)
ஆனால் ஒவ்வொரு கணமும் ஏதோ ஒரு படிநிலையில் உருவாகிறது என்ற எண்ணத்தால்
நமது கணம் என்ற \emph{கருத்தை} ஊக்குவித்தோம்
(\olref[z][story]{sec} இலுள்ள \stageshier{} ஐக் காண்க). எனவே
படிநிலைமுறைக்கு வெளியே விழும் கணங்களின் சாத்தியத்தை விலக்க விரும்புவோம்.
அதற்கேற்ப, ஒவ்வொரு கணமும் படிநிலைமுறையில் எங்காவது இடம்பெறுவதை
உறுதிசெய்யும் ஒரு புதிய அடிகோளைச் சேர்க்க வேண்டும்.

$V_\alpha$-களே நமது படிநிலைகள் என்பதால், பின்வருவதை ஓர் அடிகோளாகச்
சேர்ப்பதை வெறுமனே கருதலாம்:

\begin{defish}
\emph{ஒழுங்குமை.} $\forall A \exists \alpha\, A \subseteq V_\alpha$
\end{defish}

இது முற்றிலும் நியாயமான அணுகுமுறையாக இருக்கும். எனினும், அடுத்த பிரிவில்
விளக்கப்படும் காரணங்களுக்காக, பின்வரும் மாற்று அடிகோளை ஏற்றுக்கொள்வோம்:

\begin{axiom}[அடித்தளம்]
$(\forall A \neq \emptyset)(\exists B \in A)A \cap B = \emptyset$.
\end{axiom}

சிறிது முயற்சியுடன், $\ZFminus$ இல் அடித்தளம் ஒழுங்குமையை விளைவிக்கிறது
என்று காட்டலாம்:
\begin{defn}
ஒவ்வொரு கணம் $A$ க்கும், பின்வருமாறு கொள்க:
\begin{align*}
	\text{cl}_0(A) &= A,\\
	\text{cl}_{n+1}(A) &= \bigcup \text{cl}_n(A),\\
	\text{trcl}(A) &= \bigcup_{n < \omega} \text{cl}_{n}(A).
\end{align*}
$\text{trcl}(A)$ ஐ $A$ இன் \emph{கடப்பு மூடல்} என்று அழைக்கிறோம்.
\end{defn}
\noindent
``கடப்பு மூடல்'' என்ற பெயர் பொருத்தமானது:
\begin{prop}\ollabel{subsetoftrcl}
$A \subseteq \trcl{A}$ மற்றும் $\trcl{A}$ ஒரு கடப்புப் பண்புடைய கணம்.
\end{prop}

\begin{proof}
$A = \text{cl}_0(A) \subseteq \text{trcl}(A)$ என்பது தெளிவு.
மேலும் $x \in b \in \trcl{A}$ எனில், ஏதோ ஒரு $n$ க்கு
$b \in \text{cl}_n(A)$; ஆகவே $x \in \text{cl}_{n+1}(A)
\subseteq \text{trcl}(A)$.
\end{proof}

\begin{lem}\ollabel{lem:TransitiveWellFounded}
$A$ ஒரு கடப்புப் பண்புடைய கணம் எனில், $A \subseteq V_\alpha$ ஆக
அமையும் ஏதோ ஓர் $\alpha$ உள்ளது.
\end{lem}

\begin{proof}
\olref[ordinals][opps]{defsupstrict} இலுள்ள ``$\supstrict(X)$''
என்ற வரையறையை நினைவுகூர்ந்து, இரண்டு கணங்களை வரையறுக்கிறோம்:
\begin{align*}
	D  &= \Setabs{x \in A}{\forall \delta\ x \nsubseteq V_\delta}\\
	\alpha &= \supstrict\Setabs{\delta}{(\exists x \in A)
	(x \subseteq V_\delta \land (\forall \gamma \in \delta)x \nsubseteq V_\gamma)}
\end{align*}
$D = \emptyset$ என்க. எனவே $x \in A$ எனில், $x \subseteq
V_\delta$ ஆக அமையும் ஏதோ ஒரு $\delta$ உள்ளது; வரிசையெண்களின் நன்கு
வரிசைப்படுத்தலால் $(\forall \gamma \in \delta)x \nsubseteq
V_\gamma$ ஆகவும் தேர்ந்தெடுக்கலாம். இதனால் $\delta \in \alpha$;
எனவே \olref[spine][Valphabasic]{Valphabasicprops} இன்படி
$x \in V_\alpha$. ஆகவே தேவைப்பட்டவாறு $A \subseteq V_{\alpha}$.

எனவே $D = \emptyset$ என்று காட்டினால் போதும். முரண்பாட்டிற்காக,
அவ்வாறு இல்லை என்க. அடித்தளத்தின்படி, $D \cap B = \emptyset$ ஆக
அமையும் ஏதோ ஒரு $B \in D \subseteq A$ உள்ளது. $x \in B$ எனில்,
$A$ கடப்புப் பண்புடையதால் $x \in A$; மேலும் $x \notin D$ என்பதால்
$\exists \delta\ x \subseteq V_\delta$. இப்போது
\[
	\beta = \supstrict\Setabs{\delta}{(\exists x \in B)(x \subseteq V_\delta \land (\forall \gamma < \delta)x \nsubseteq V_\gamma)}.
\]
முன்புபோலவே, $B \subseteq V_\beta$; இது $B \in D$ என்ற கூற்றுடன்
முரண்படுகிறது.
% SOURCE CORRECTION TA-STH-015: restore uppercase B in the defining set for beta.
\end{proof}

\begin{thm}\ollabel{zfentailsregularity}
ஒழுங்குமை நிறைவேறுகிறது.
\end{thm}

\begin{proof}
$A$ ஐ நிலைப்படுத்துக; \olref{subsetoftrcl} இன்படி $A \subseteq
\trcl{A}$, மேலும் இது கடப்புப் பண்புடையது. எனவே
\olref{lem:TransitiveWellFounded} இன்படி ஏதோ ஒரு $\alpha$ க்கு
$A \subseteq \trcl{A} \subseteq V_\alpha$.
\end{proof}

இம்முடிவுகள் $\ZFminus$ ஆனது
$\emph{அடித்தளம்}\Rightarrow\emph{ஒழுங்குமை}$ என்ற நிபந்தனைவாக்கியத்தை
நிறுவுகிறது என்று காட்டுகின்றன. \olref[spine][rank]{zfminusregularityfoundation}
இல், $\ZFminus$ ஆனது $\emph{ஒழுங்குமை}\Rightarrow\emph{அடித்தளம்}$
என்ற நிபந்தனைவாக்கியத்தை நிறுவுகிறது என்றும் காட்டுவோம். எனவே அடித்தளமும்
ஒழுங்குமையும் ($\ZFminus$ தொடர்பில்) \emph{சமானவை}. ஆனால் இதன் பொருள்,
$\ZFminus$ கொடுக்கப்பட்டால், அடித்தளம் ஒழுங்குமைக்குச் சமானமானது என்று
குறிப்பிட்டு அடித்தளத்தை நியாயப்படுத்தலாம். மேலும் \stageshier{} இன்
அடிப்படையில் ஒழுங்குமையை உடனடியாக நியாயப்படுத்தலாம்.

\end{document}
